%%% CSE 4316 Project Charter - WORKED EXAMPLE (Full-Stack Web App). Follows ISO/IEC/IEEE 16326:2019.
\documentclass[11pt,letterpaper]{article}
\usepackage[margin=1.0in]{geometry}
\usepackage[T1]{fontenc}\usepackage[bitstream-charter]{mathdesign}\usepackage[latin1]{inputenc}
\usepackage{amsmath}\usepackage{xcolor}\usepackage{cite}\usepackage{hyphenat}
\usepackage{graphicx}\usepackage{float}\usepackage{subfigure}\usepackage{sectsty}
\usepackage[compact]{titlesec}\usepackage[tablegrid]{vhistory}\usepackage{longtable}
\allsectionsfont{\color{accentcolor}\scshape\selectfont}
\definecolor{accentcolor}{rgb}{0.0,0.0,0.5}
\newcommand{\teamname}{Team Fullstack}
\newcommand{\productname}{CivicPulse Community Service-Request Dashboard}
\newcommand{\coursename}{CSE 4316: Senior Design I}
\newcommand{\semester}{Fall 2026}
\newcommand{\docname}{Project Charter}
\newcommand{\department}{Department of Computer Science \& Engineering}
\newcommand{\university}{The University of Texas at Arlington}
\newcommand{\authors}{Alan Turing \\ Grace Hopper \\ John Von Neumann \\ Ada Lovelace \\ Charles Babbage}
\usepackage{fancyhdr}\pagestyle{fancy}\usepackage{lastpage}
\lhead{}\chead{}\rhead{}\lfoot{\footnotesize \teamname \ - \semester}\cfoot{}
\rfoot{\footnotesize page \thepage\ of \pageref{LastPage}}
\renewcommand{\headrulewidth}{0.0pt}\renewcommand{\footrulewidth}{0.4pt}
\usepackage[runin]{abstract}\abslabeldelim{\quad}\setlength{\abstitleskip}{-10pt}
\renewcommand{\abstractname}{}\renewcommand{\abstracttextfont}{\color{accentcolor} \small \slshape}

\begin{document}
{\centering \huge \color{accentcolor} \sc \textbf{\department \\ \university} \par}
\vspace{1 in}
{\centering \huge \color{accentcolor} \sc \textbf{\docname \\ \coursename \\ \semester} \par}
\vspace{0.5 in}
\begin{figure}[h!]\centering\includegraphics[width=0.60\textwidth]{images/test_image}\end{figure}
\vspace{0.5 in}
{\centering \huge \color{accentcolor} \sc \textbf{\teamname \\ \productname} \par}
\vspace{0.5 in}{\centering \large \sc \textbf{\authors} \par}
\newpage
\begin{versionhistory}
  \vhEntry{0.1}{09.02.2026}{GH}{document creation}
  \vhEntry{1.0}{09.15.2026}{AT|GH|CB}{official release}
\end{versionhistory}
\newpage\tableofcontents\newpage

\section{Problem Statement}
Residents report non-emergency issues --- potholes, broken streetlights, illegal
dumping --- through phone calls and email that land in scattered inboxes. Public
works staff cannot see the full backlog, duplicates pile up, and residents get no
status back, eroding trust. The city needs a single place for residents to report
issues and for staff to see, prioritize, and resolve them, with visible status.

\section{Methodology}
Team Fullstack will build \productname, a web application with a public resident
portal for submitting and tracking geolocated requests and a staff dashboard for
triage, assignment, and status updates. In short: we are building a website where
residents report neighborhood issues and city staff manage them to resolution on a
shared dashboard.

\section{Value Proposition}
For the city, CivicPulse consolidates intake, removes duplicate handling, and gives
residents self-service status, reducing call volume and improving accountability.
For the university, the project exercises a complete secure multi-tier web
application --- authentication, role-based access, geospatial data, and reporting.

\section{Development Milestones}
\begin{itemize}
  \item Project Charter first draft - September 2026
  \item System Requirements Specification - October 2026
  \item Architectural Design Specification - November 2026
  \item Demonstration of request submission \& map view - December 2026
  \item Detailed Design Specification - February 2027
  \item Demonstration of staff triage \& assignment - March 2027
  \item Demonstration of status notifications - April 2027
  \item CoE Innovation Day poster presentation - April 2027
  \item Final Project Demonstration (end-to-end workflow) - May 2027
\end{itemize}
\newpage

\section{Background}
The department tracks requests in spreadsheets and email, with no resident-facing
status. They asked the CSE program for a modern web front door integrated with their
existing single sign-on. No prior relationship; requirements from two interviews and
a review of their current intake spreadsheet.

\section{Related Work}
Commercial 311/CRM platforms (e.g. SeeClickFix-style products) exist \cite{Rubin2012}
but are subscription-priced and not tailored to the city's SSO and categories.
CivicPulse targets a lightweight, self-hostable alternative. (Add $\geq$5 IEEE
references in the final version.)

\section{System Overview}
A responsive web frontend (resident portal + staff dashboard) talks to a backend API
that enforces authentication and role-based access, persists requests with location
in a database, and sends status notifications by email. A map component visualizes
open requests. See the ADS for the three-tier design.

\section{Project Management Approach}
Iterative Scrum \cite{Rubin2012, SchwaberSutherland2020} per ISO/IEC/IEEE 16326:2019
\cite{ISO16326} within ISO/IEC/IEEE 12207:2017 \cite{ISO12207}. Two-week sprints;
estimation in task-hours; progress by burndown and by requirement/test coverage.
Trunk-based Git with CI; tagged document releases.

\section{Roles \& Responsibilities}
Sponsor contact: public-works IT liaison. Product Owner: Grace Hopper. Scrum Master:
Alan Turing. Frontend: Ada Lovelace; Backend/API: John Von Neumann; Data/geo:
Charles Babbage; Auth/notifications: Grace Hopper; Integration/test: Alan Turing.

\section{Cost Proposal}
Mainly cloud hosting for a staging environment plus a mapping/geocoding API (free
tier where possible); no custom hardware. Funded by the default CSE stipend; largest
line item is hosting.

\section{Facilities \& Equipment}
Senior Design lab workstations; a cloud project for staging; access to a sponsor
test SSO realm by the 3rd sprint.

\section{Assumptions}
\begin{itemize}
  \item A sponsor SSO/OIDC test realm will be available by the 3rd sprint.
  \item The city will provide the list of request categories and routing rules.
  \item A free-tier geocoding/mapping API will meet demo volume.
  \item Residents access the portal from modern mobile and desktop browsers.
  \item Email delivery via a provided SMTP relay.
\end{itemize}

\section{Constraints}
\begin{itemize}
  \item Final demonstration must be completed by May 1st, 2027.
  \item Total development cost must not exceed \$800.
  \item The application must meet WCAG 2.2 AA accessibility.
  \item Resident personal data handling must follow the city's privacy policy.
  \item Staff functions must be restricted by role (no public access).
\end{itemize}

\section{Risks}
\begin{table}[h]\resizebox{\textwidth}{!}{
\begin{tabular}{|l|l|l|l|}
\hline
\textbf{Risk description} & \textbf{Probability} & \textbf{Loss (days)} & \textbf{Exposure (days)} \\ \hline
SSO/OIDC integration with city realm delayed & 0.40 & 12 & 4.8 \\ \hline
Scope creep in staff dashboard features & 0.45 & 10 & 4.5 \\ \hline
Geocoding API limits/costs on real volume & 0.30 & 8 & 2.4 \\ \hline
Accessibility rework late in the project & 0.35 & 7 & 2.5 \\ \hline
Email deliverability issues & 0.20 & 6 & 1.2 \\ \hline
\end{tabular}}
\caption{Highest-exposure project risks}
\end{table}

\section{Documentation \& Reporting}
Documentation deliverables in Git, released at their planning sprints; recurring
sprint artifacts and engineering notebooks each sprint. Closeout: deployed
application, poster, public web page, demo video, documented source (MIT),
infrastructure-as-code deployment scripts, and admin/user manuals.

\bibliographystyle{reference/IEEEtran_custom}
\bibliography{reference/refs}{}
\end{document}
